\section{Threats to Validity}
\label{sec:threats}

\subsection{Construct Validity}

The largest construct threat is the relationship between the available binary
label and the stated ransomware domain. The repository does not include original
file identities or a label-definition document proving that every positive row
is ransomware. Calling the measured task ransomware detection would therefore
overstate what the data establish. We describe it as binary classification of a
local PE-feature table and treat ransomware as the motivating application.

The acquisition signals also require precise interpretation. $H$ is entropy of
positive versus negative SHAP mass, not predictive class entropy. $V$ is variance
over signed coordinates for one explanation, not uncertainty across trees or
models. The labels ``directional conflict'' and ``dispersion'' are used to avoid
claiming that either quantity is a calibrated epistemic uncertainty measure.

Fidelity is similarly multi-dimensional. A large deletion effect supports the
relationship between a ranking and fitted-model behavior under a defined
intervention. It does not prove causality, human usefulness, stability under all
perturbations, or semantic correctness. EC and CF are implementation-specific
geometric statistics. We avoid collapsing them into OFS because a composite
would obscure rather than solve these construct differences.

\subsection{Internal Validity}

The fixed row split prevents direct overlap among train, validation, and test
indices, and the audit found no exact duplicate feature vectors. It does not rule
out near-duplicates, related binaries, or shared acquisition sources across
partitions. Such dependence could inflate performance. Family-, source-, or
time-aware splitting cannot be performed without metadata.

Active-learning methods share initial sets within seeds, budgets, hyperparameters,
and test data. This controls important confounders, but fixed acquisition weights
were not selected in an independent outer loop. The policy comparison may depend
on those choices. The oracle uses existing labels and has no annotation error;
real analysts could change both selected-label quality and measured cost.

The fidelity model is trained separately from the active-learning models. This
is intentional for a fixed-model diagnostic, but it prevents a claim that the AL
improvement caused the fidelity outcome. Training means are calculated only from
the training partition, which avoids test leakage. Nonetheless, mean replacement
can induce out-of-distribution vectors, and top-SHAP versus random masks may have
different numeric magnitudes despite equal selected-coordinate counts.

No internal-validity claim is made for PID performance. The runner rejects
missing, simulated, or untraceable dynamic evidence, and the default experiment
did not execute fusion. This avoids one confound---constructing $\pdynamic$ from
the target label---but leaves RQ3 empirically unanswered.

\subsection{Statistical Conclusion Validity}

Five active-learning seeds share one test set. Their sample SD quantifies
sensitivity to initial labels and query sequences, not uncertainty over the
population of possible malware samples. Paired seed differences are appropriate
for the controlled comparison, but five pairs provide limited power and do not
justify strong distributional assumptions. We report all paired signs and avoid
a fragile significance claim.

FPR comparisons depend on thresholds selected from validation at a target TPR.
Because achieved test TPR varies, a lower FPR is not automatically a uniformly
better detector. Both rates are reported. AP is calculated from ranked
predictions using the standard step-wise definition and is not interchangeable
with trapezoidal PR area \cite{saito2015precision}.

The fidelity bootstrap resamples test rows while retaining paired original and
masked scores. Its interval is conditional on one trained model. Twenty random
masks characterize mask variability for that model and test set, not training or
dataset variability. EC uses 79,800 pairwise entries, but these entries share 400
samples, so they cannot be treated as independent observations for a conventional
Pearson test.

\subsection{External Validity}

The pilot is small relative to contemporary malware corpora and has no verified
upstream provenance. Results may change across collection periods, organizations,
malware families, benign software distributions, packers, and feature-extractor
versions. The study is neither a temporal holdout nor a family holdout, and it
does not establish zero-day performance. External validation should use a
manifested dataset with file hashes, collection dates, family/source metadata,
licensing, and independently documented labels.

The models use EMBER-style static coordinates and LightGBM. The acquisition
signals can be calculated for other additive explanations, but their behavior and
cost may differ for raw-byte networks, graph models, or sequence models. Likewise,
mean masking is not directly transferable to categorical, sparse, or structured
representations without a domain-specific baseline.

\subsection{Operational and Reproducibility Validity}

Recorded timings were obtained in a small, deterministic two-thread experiment.
They exclude several deployment costs and cannot establish endpoint throughput or
SOC response time. No real sandbox trace is included, so the work does not show
resistance to sleep evasion or environmental checks.

The evidence package provides hashes, versions, configurations, checkpoints,
queries, predictions, masks, and successful replay counts. This supports
computational reproducibility for the retained run. The compact review package
does not commit all large checkpoints and per-sample artifacts; a fresh full run
is required to regenerate them. Moreover, source and output hashes verify file
identity rather than the truth of external metadata. Reproducibility and data
authenticity are related but distinct properties.
